\documentclass[11pt]{article}
\usepackage[utf8]{inputenc}
\usepackage{amsmath,amssymb,amsthm}
\usepackage[margin=1in]{geometry}
\usepackage{hyperref}
\usepackage{xcolor}
\usepackage{enumitem}
\usepackage{newunicodechar}
\newunicodechar{ℝ}{\ensuremath{\mathbb{R}}}
\newunicodechar{⟨}{\ensuremath{\langle}}
\newunicodechar{⟩}{\ensuremath{\rangle}}
\newunicodechar{Γ}{\ensuremath{\Gamma}}
\newunicodechar{·}{\ensuremath{\cdot}}
\newunicodechar{²}{\ensuremath{^2}}
\newunicodechar{φ}{\ensuremath{\varphi}}
\newunicodechar{ψ}{\ensuremath{\psi}}
\newunicodechar{λ}{\ensuremath{\lambda}}
\newunicodechar{≠}{\ensuremath{\ne}}
\newunicodechar{≤}{\ensuremath{\le}}

\setlength{\emergencystretch}{3em}

\newcommand{\userbox}[1]{\par\medskip\begin{quote}\small\textbf{\textcolor{blue!60}{DM:}} #1\end{quote}\medskip}
\newcommand{\claudebox}[1]{\par\medskip\begin{quote}\small\textbf{\textcolor{orange!60!black}{Claude:}} #1\end{quote}\medskip}
\newcommand{\gptbox}[1]{\par\medskip\begin{quote}\small\textbf{\textcolor{green!50!black}{GPT-5.5 Pro:}} #1\end{quote}\medskip}

\title{Tide retrospective:\\\emph{Separable-potential factorisation infrastructure}}
\author{Daniel Murfet (with Claude Opus 4.7)}
\date{6 May 2026}

\begin{document}
\maketitle

\begin{abstract}
A single Tide-loop excursion in the \texttt{laplace} seabed, exposing
the additively-separable factorisation pattern that the existing
\texttt{SemiDegenerate.lean} (harmonic + quartic) and
\texttt{PureQuartic.lean} (quartic + quartic) both re-derive case by
case. The new module\footnote{Module:
\texttt{Laplace.TwoD.AddSeparable}.} adds 156 lines of Lean to the
seabed, with five generic theorems: a definition, the Boltzmann
decomposition, partition-function factorisation, separable-observable
factorisation, separable-Gibbs-expectation factorisation, and the
mixed-covariance vanishing $\mathrm{Cov}[f \otimes 1,\, 1 \otimes g] =
0$. Zero \texttt{sorry}, zero \texttt{axiom}, zero
\texttt{native\_decide}; clean diagnostics on first compile. The
interesting moves were upstream: GPT-5.5 Pro flagged a wrong covariance
formula in the candidate set (the cross-mean terms
$(a_1 b_2 + a_2 b_1) \cdot \langle x \rangle_U \cdot \langle y \rangle_V$
actually cancel under separable Gibbs, never appear in the final
covariance), and recommended a tighter set of hypotheses on the
partition-factorisation that the Bochner unconditional form would not
have caught.
\end{abstract}

\section{Setting}

This is the first ``infrastructure abstraction'' Tide step in the
\texttt{laplace} seabed: it does not extend the formalised territory
outward by adding a new headline theorem, but factors the existing
shoreline by extracting a shared template from two parallel files. Both
\texttt{SemiDegenerate.lean} (harmonic + quartic, 514 lines) and
\texttt{PureQuartic.lean} (quartic + quartic, 478 lines) carry a
structurally identical proof skeleton: define a 2D potential as a sum,
prove the Boltzmann factor decomposes as a product, apply Mathlib's
\texttt{integral\_prod\_mul} to factor the partition function, then
expand to closed form via the 1D pieces.

The Tide step abstracts that template once for any
$L(x, y) = U(x) + V(y)$. Future 2D potentials (notably the
quartic--sextic mixed and any subsequent monomial--monomial pair)
specialise rather than copy-paste-rewrite. The deferred-but-cost-positive
refactor was identified at the end of the 2D pure-quartic
retrospective\footnote{\texttt{lean/laplace/retrospectives/2026-05-02-tide-2d-pure-quartic.tex},
\textsection Lessons.} as the next abstraction worth lifting; the 6 May
candidates survey lifted it to candidate A3 and recommended it as the
high-yield refactor pick of the post-Tide-10 batch. The user picked
A3 directly from the survey digest.

The seabed was \texttt{main} at \texttt{1e3802a} (the merge of Tide 11,
\texttt{harmonic-gibbsregularity}); the new branch is
\texttt{tide/separable-potential}, off \texttt{main} (no chain needed
this time --- all prior tides are merged).

\section{The thing formalised}

For any pair of potentials $U, V : \mathbb{R} \to \mathbb{R}$, the
additively-separable 2D potential is $L(x, y) = U(x) + V(y)$. The
module exposes five theorems, of which the headline two are:

\begin{quote}\itshape
Under integrability of each marginal Boltzmann weight,
\[
  Z_{2D}(t) \;=\; Z_U(t) \cdot Z_V(t)
\]
where $Z_{2D}(t) = \int e^{-tL(x,y)}\,dx\,dy$ and $Z_U(t) = \int
e^{-tU(x)}\,dx$, $Z_V(t) = \int e^{-tV(y)}\,dy$. For separable
observables $\varphi(x, y) = f(x)$, $\psi(x, y) = g(y)$, when both
1D partition functions are nonzero,
\[
  \mathrm{Cov}_{2D, t}[\varphi, \psi] \;=\; 0.
\]
\end{quote}

\noindent The Lean signatures (abbreviated; full file at
\texttt{Laplace/TwoD/AddSeparable.lean}):

\begin{verbatim}
noncomputable def addSeparable (U V : ℝ → ℝ) : ℝ × ℝ → ℝ :=
  fun z ↦ U z.1 + V z.2

theorem partitionFunction_addSeparable_factor
    {U V : ℝ → ℝ} {t : ℝ}
    (hU : Integrable (fun x ↦ Real.exp (-(t * U x))))
    (hV : Integrable (fun y ↦ Real.exp (-(t * V y)))) :
    partitionFunction (addSeparable U V) t =
      Laplace.partitionFunction U t * Laplace.partitionFunction V t

theorem integral_separable_addSeparable
    {U V : ℝ → ℝ} {t : ℝ} (f g : ℝ → ℝ)
    (hf : Integrable (fun x ↦ f x * Real.exp (-(t * U x))))
    (hg : Integrable (fun y ↦ g y * Real.exp (-(t * V y)))) :
    (∫ z : ℝ × ℝ, f z.1 * g z.2 * Real.exp (-(t * addSeparable U V z))) =
      (∫ x, f x * Real.exp (-(t * U x))) * (∫ y, g y * Real.exp (-(t * V y)))

theorem gibbsCov_addSeparable_fst_snd_eq_zero
    {U V : ℝ → ℝ} {t : ℝ} (f g : ℝ → ℝ)
    (hZU_ne : Laplace.partitionFunction U t ≠ 0)
    (hZV_ne : Laplace.partitionFunction V t ≠ 0)
    (hU : Integrable ...)  (hV : Integrable ...)
    (hf : Integrable ...)  (hg : Integrable ...) :
    gibbsCov (addSeparable U V) t (fun z ↦ f z.1) (fun z ↦ g z.2) = 0
\end{verbatim}

The mixed-covariance vanishing is the structural payoff of separability:
under a product Gibbs measure, the covariance between an observable
that depends only on the first coordinate and one that depends only on
the second is identically zero. This is the lemma that the existing
\texttt{cov\_affine\_*} theorems implicitly use when they observe that
the $a_1 b_2$ and $a_2 b_1$ cross terms drop out.

Numerical sanity: not feasible (structural identities, not specific
numerical values). The seabed's existing 2D files have already
validated the concrete instances numerically.

\section{Proof strategy}

All four non-trivial theorems pivot on a single Mathlib lemma:

\begin{quote}\itshape
$\int_{\mathbb{R}^2} f(z_1)\, g(z_2)\, d(\mu \times \nu)
\;=\; \big(\int f\, d\mu\big) \cdot \big(\int g\, d\nu\big)$.
\end{quote}

This is \texttt{MeasureTheory.integral\_prod\_mul} in Mathlib, and it
is unconditional: when either factor is non-integrable the equation
holds with both sides equal to $0$ by Bochner convention. After the
Boltzmann decomposition $e^{-t(U(z_1) + V(z_2))} = e^{-t U(z_1)} \cdot
e^{-t V(z_2)}$ via \texttt{Real.exp\_add}, all of partition,
integral-of-separable-observable, and Gibbs-expectation factorisations
are one application of \texttt{integral\_prod\_mul} away.

The mixed-covariance corollary is the only real proof in the file:
$\mathrm{Cov}[\varphi, \psi] = \langle \varphi \psi \rangle - \langle
\varphi \rangle \langle \psi \rangle$. The first term factors via the
expectation lemma at $(f, g)$. The second term factors as $\langle f
\rangle_U \cdot 1 \cdot 1 \cdot \langle g \rangle_V$, after applying the
expectation lemma to each of $\varphi = f \otimes 1$ and $\psi = 1
\otimes g$ separately and using
\texttt{Laplace.gibbsExpectation\_const} to collapse the constant
factors. Then the two products are equal and \texttt{ring} closes.

The proof of the constant-factor collapse is a four-line $\beta$-eta
gymnastics: the helper expects an observable of the form $f(z) = h(z_1) \cdot
k(z_2)$, and we want to apply it at $k = 1$ to get $\varphi = f \otimes 1
= f(z_1) \cdot 1$, then rewrite $f(z_1) \cdot 1 = f(z_1)$. The
\texttt{simpa using hV} pattern handles the integrability witness for
$1 \cdot e^{-t V(y)}$ via \texttt{one\_mul}.

\section{Roadblocks and resolutions}

The deliberation surfaced two issues; the formalisation surfaced one.

\paragraph{GPT correction on Candidate B's covariance formula.}
My v1 statement of the affine-covariance template (Candidate B) had:
\[
  \mathrm{Cov}_t[a_1 x + a_2 y + c,\, b_1 x + b_2 y + d]
  = a_1 b_1 \mathrm{Var}_U + a_2 b_2 \mathrm{Var}_V
  + (a_1 b_2 + a_2 b_1) \langle x \rangle_U \langle y \rangle_V.
\]
GPT pointed out that under a separable Gibbs law, the cross terms
$\langle x \rangle_U \langle y \rangle_V$ \emph{cancel} between the
joint expectation $\langle (a_1 x + a_2 y + c)(b_1 x + b_2 y + d)
\rangle$ and the product of marginals $\langle a_1 x + a_2 y + c
\rangle \cdot \langle b_1 x + b_2 y + d \rangle$. The clean formula is
just
\[
  \mathrm{Cov}_t[\varphi, \psi]
  = a_1 b_1 \cdot \mathrm{Cov}_U[x, x] + a_2 b_2 \cdot \mathrm{Cov}_V[y, y].
\]
This wasn't load-bearing for the present tide (we deferred B), but it
matters for the follow-up tide that picks up the affine template. The
existing files' \texttt{cov\_affine\_*} theorems do the cancellation
inline in their algebraic-manipulation step; the abstracted version
states the clean formula directly.

\paragraph{Hypothesis form on `partitionFunction_addSeparable_factor`.}
My v1 had the partition factorisation \emph{unconditional}, relying on
the Bochner convention that $\int = 0$ on non-integrable functions
makes the equation true even without integrability. GPT recommended
adding integrability hypotheses anyway, on the grounds that downstream
files always have integrability witnesses available and the
hypothesis-bearing form is the one the reader naturally chains
through:

\gptbox{For Bochner $\int$, unconditional factorisation statements
are brittle because the theorem you will actually use is typically
the integrable version. \ldots~Downstream files will naturally be
able to apply [the hypothesis-bearing form], since for concrete
potentials they already prove finiteness/integrability of the
partition function.}

The committed file takes the integrability hypotheses but doesn't use
them in the proof body, since \texttt{integral\_prod\_mul} is
unconditional. The hypotheses are documentation / API discipline rather
than load-bearing. This is unusual but defensible: a future Mathlib
refactor that adds integrability hypotheses to \texttt{integral\_prod\_mul}
would not break our callers.

\paragraph{Concurrent-agent branch race (recurrence).}
A parallel session created its own tide branch off \texttt{main}
while my work was in flight, and \texttt{git commit} fired against
\emph{that} branch's checkout instead of the intended
\texttt{tide/separable-potential}. The fix was a two-line
cleanup\footnote{\texttt{git branch -f tide/separable-potential <SHA>}
to relocate the commit, then \texttt{git branch -f <other-tide> main}
to restore the parallel session's branch baseline.}. This is the
third such recurrence (2D pure-quartic, sextic-moments, now
separable-potential). The structural fix --- run each Tide agent in
its own \texttt{git worktree} --- remains a recommended skill change
rather than landed infrastructure. (After the cleanup, this Tide
moved to a worktree at \texttt{lean/laplace-separable/} for the
remainder of the run; the parallel-session contention disappeared
immediately.)

Beyond these three, no proof-side roadblocks. All five theorems
typechecked clean on the first compile.

\section{What was Mathlib, what was new}

\paragraph{Mathlib provided.} \texttt{MeasureTheory.integral\_prod\_mul}
was the load-bearing lemma; the four non-trivial theorems all reduce to
one application of it after a \texttt{Real.exp\_add} for the Boltzmann
decomposition. The 1D Gibbs framework
(\texttt{Laplace.partitionFunction}, \texttt{Laplace.gibbsExpectation},
\texttt{Laplace.gibbsCov}, \texttt{Laplace.gibbsExpectation\_const})
came from the seabed's own \texttt{Laplace.Gibbs} module. The 2D Gibbs
framework (the \texttt{Laplace.TwoD.}-prefixed wrappers used in
\texttt{SemiDegenerate.lean}) came from the same.

\paragraph{What was new in this tide.} The five theorems in the new
module: the additively-separable potential definition, its Boltzmann
decomposition lemma, partition-function factorisation,
separable-observable integral factorisation,
Gibbs-expectation-of-separable factorisation, and the mixed-covariance
vanishing. None of the five theorems introduces new mathematical
content; the value is in stating the abstraction once with the right
hypothesis package, where ``the right'' was determined by GPT-5.5 Pro's
critique of the v1 hypotheses.

\section{Lessons}

\begin{itemize}
\item \emph{The deliberation step pays for itself when the candidate
is structural.} For algebraic / asymptotic Tides (kappa3-harmonic,
anharmonic-kappa3) GPT's contribution is usually a strategic move
inside the proof. For an abstraction Tide like this one, GPT's
contribution is on the \emph{statement} --- specifically, the
hypothesis package and the algebraic cancellation that hides in the
naive form. Both interventions saved follow-up cycles in the
formalisation.

\item \emph{Unconditional factorisation theorems are an anti-pattern
even when valid.} \texttt{integral\_prod\_mul} is unconditional in
Mathlib, but stating our partition factorisation unconditionally
would have invited downstream confusion. Take integrability
hypotheses anyway, even if the proof body doesn't use them; they
document the API and survive future Mathlib tightening.

\item \emph{Concurrent Tide agents continue to need filesystem
isolation.} Three tides have now been bitten by the
shared-checkout branch race. This recurrence makes the case for
landing \texttt{git worktree} support in the Tide skill's preflight
check rather than continuing to hand-craft the recovery dance.
\end{itemize}

\section{Follow-ups}

\begin{itemize}
\item \emph{Refactor the existing 2D files to use the abstraction.}
Both \texttt{SemiDegenerate.lean} and \texttt{PureQuartic.lean} carry
\texttt{partitionFunction\_factor} and \texttt{integral\_pow\_pow\_factor}
lemmas that are now corollaries of the new module's theorems at
concrete $(U, V)$ choices. The follow-up tide replaces $\sim$30 lines
per file with $\sim$5-line specialisations.

\item \emph{Generic Gibbs-covariance algebra.} GPT recommended a small
covariance API in the base Gibbs files: \texttt{gibbsCov\_symm},
\texttt{gibbsCov\_add\_left/right}, \texttt{gibbsCov\_const\_left/right},
\texttt{gibbsCov\_smul\_left/right}. Once those land, the affine
covariance template is short and pleasant, and Candidate B from this
deliberation becomes essentially trivial.

\item \emph{Mixed-power 2D moments (A4 from the 6 May survey).} Now a
one-line corollary of \texttt{integral\_separable\_addSeparable} at
$f = x \mapsto x^m$, $g = y \mapsto y^n$. Pairs naturally with the
``even-monomial template'' tide (A1).

\item \emph{Quartic--sextic mixed potential (A2 from the 6 May survey).}
Now a specialisation of the new abstraction with closed-form
$U = $ quartic and $V = $ sextic. The deferred 2D Tide.

\item \emph{Affine-covariance template (Candidate B / future tide).}
Picking up the deferred piece, ideally after the generic Gibbs-covariance
algebra lands. The clean statement (no surviving cross-mean terms) is
in this retrospective for the future tide to consume.
\end{itemize}

\section*{Acknowledgements}

GPT-5.5 Pro contributed both deliberation-stage corrections that
shaped the final theorem statements: the Bochner-integrability
hypothesis discipline and the algebraic cancellation in the affine
covariance formula. The full consult and the tide log are preserved in
the primer project's tide-log directory under the SRI repo.

\end{document}
